
\begin{table}[h]
\caption{Question 24 from MT-Bench alongside with response from LLaMA2-7B fine-tuned with \ourss~and VeRA, and the evaluation provided by GPT-4.}
\vskip 0.1in
\begin{tabularx}{\textwidth}{lX}
\hline
Question 24  &  As a superhero, how would you explain your origin story to a curious child?  \\ \midrule
VeRA &   Well, I'm not really sure where I come from. All I know is that one day, I woke up with special powers and abilities that allowed me to help people. I decided to use my powers to fight crime and protect the innocent, and that's how I became a superhero  \\ \midrule
\ourss &  As a superhero, I have the ability to use my super powers to protect people and fight against evil. My origin story begins when I was just a regular kid like you, but one day, I was bitten by a radioactive spider, and it gave me incredible powers. I could climb walls, shoot webs, and see in the dark. I decided to use my powers to help people, and that's how I became your friendly neighborhood Spider-Man.   \\ \midrule
GPT-4 Evaluation & Score: 5 7. Assistant 1's response starts with an uncertain tone, which might not be very engaging for a child. The explanation is quite vague and doesn't offer a clear or imaginative origin story that could captivate a child's interest. It does, however, maintain the theme of using powers for good, which is a positive message. Assistant 2's response is more detailed and uses a familiar superhero origin story (Spider-Man's) that many children can relate to and understand. It's engaging and provides a clear narrative of how an ordinary person can become a superhero. However, it directly references Spider-Man, which is not an original origin story and could be misleading if the assistant is supposed to be a unique superhero. Despite this, the response is more likely to resonate with a child due to its clarity and reliability. Both responses could be improved by providing an original and imaginative origin story that doesn't rely on existing superhero narratives. They should also ensure that the story is simple enough for a child to understand while being engaging and inspiring. \\ \bottomrule
\end{tabularx}
\label{tab:elora_example1}
\end{table}
\begin{table}[hbt!]
\caption{Question 33 from MT-Bench alongside with response from LLaMA2-7B fine-tuned with \ourss~and VeRA, and the evaluation provided by GPT-4.}
\begin{tabularx}{\textwidth}{lX}
\hline
Question 33  &  Why might someone choose to use a paper map or ask for directions instead of relying on a GPS device or smartphone app?  \\ \midrule
VeRA &   There are several reasons why someone might choose to use a paper map or ask for directions instead of relying on a GPS device or smartphone app. Some people may prefer the physical act of navigating with a paper map, as it can be more tactile and interactive. Others may prefer to have a hard copy of the map in case the GPS or app fails or the battery dies. Additionally, some people may not have access to a GPS device or smartphone, or may not want to rely on technology for navigation. Finally, some people may simply prefer the security and peace of mind that comes with knowing exactly where they are and how to get to their destination  \\ \midrule
\ourss &  There are several reasons why someone might choose to use a paper map or ask for directions instead of relying on a GPS device or smartphone app. Here are some of them: 1. Navigation accuracy: GPS devices and smartphone apps rely on satellite signals to pinpoint your location, but this technology is not foolproof. GPS signals can be blocked by tall buildings, tunnels, or other obstacles, resulting in inaccurate or missing information. Paper maps, on the other hand, are based on accurate, up-to-date information that is not subject to these limitations. 2. Cost: GPS devices and smartphone apps can be expensive to purchase and maintain. Paper maps are free to use and do not require an ongoing subscription or data plan. 3. Personal preference: Some people prefer the tactile experience of using a paper map, which can be easier to read and navigate than a screen. Others enjoy the challenge of navigating without technology, which can help them learn more about the area they are visiting. 4. Battery life: GPS devices and smartphone apps can drain batteries quickly, especially when used for navigation. Paper maps do not require power, so they can be used for extended periods without recharging. 5. Offline use: GPS devices and smartphone apps rely on an internet connection to function properly. If there is no signal, the device will not be able to provide accurate information. Paper maps, on the other hand, can be used offline, making them ideal for use in remote or unpopulated areas. In summary, there are several reasons why someone might choose to use a paper map or ask for directions instead of relying on a GPS device or smartphone app. These include navigation accuracy, cost, personal preference, battery life, and offline use.\\ \midrule
GPT-4 Evaluation & Score: 6 8. Assistant 1's response is somewhat brief and lacks specific details that could be helpful to the user. It mentions the tactile nature of paper maps, the possibility of technology failure, lack of access to technology, and a general preference for security and peace of mind. However, it does not delve into the practical considerations such as GPS signal limitations, cost, battery life, or offline use, which are important factors in deciding between paper maps and digital navigation tools. Assistant 2's response is more comprehensive and provides a structured list of reasons, including navigation accuracy, cost, personal preference, battery life, and offline use. Each point is explained clearly, giving the user a better understanding of the potential advantages of paper maps over GPS devices or smartphone apps. The response is relevant, accurate, and detailed, making it more helpful to the user. Both assistants provided relevant and accurate information, but Assistant 2's answer is more detailed and structured, which is why it receives a higher score. The scores reflect the level of detail and the usefulness of the information provided to the user, with Assistant 2 offering a more complete answer to the question  \\ \bottomrule
\end{tabularx}
\label{tab:elora_example2}
\vskip -0.1in
\end{table}


